% Part: incompleteness
% Chapter: introduction
% Section: definitions

\documentclass[../../../include/open-logic-section]{subfiles}

\begin{document}

\olfileid{inc}{int}{def}

\olsection{تعریف‌ها}

برای اجرای پروژهٔ هیلبرت در صوری‌سازی ریاضیات و نشان‌دادن اینکه چنین
صوری‌سازی‌ای سازگار و کامل است، نخستین کار انتخاب زبان، چارچوب منطقی
و دستگاهی از اصول موضوعه است. برای مقصود ما فرض کنیم ریاضیات را
می‌توان در زبانی مرتبه‌اول صوری کرد؛ یعنی مجموعه‌ای از !!{constant}s،
!!{function}s و !!{predicate}s وجود دارد که همراه با ادات و سورهای منطق
مرتبه‌اول امکان بیان ادعاهای ریاضی را می‌دهد. بیشتر مردم موافق‌اند
چنین زبانی وجود دارد: زبان نظریهٔ مجموعه‌ها که در آن~$\in$ تنها نماد
غیرمنطقی است. اینکه چنین زبان ساده‌ای این‌همه بیانگر باشد، البته در
نگاه نخست ادعایی بسیار نامحتمل است و اثبات اینکه عملاً همهٔ ریاضیات
را می‌توان با این واژگان بسیار کمینه بیان کرد به کار فراوانی نیاز
داشت. برای سادگی، فعلاً بحث را به حساب، یعنی بخشی از ریاضیات که فقط
با اعداد طبیعی~$\Nat$ سروکار دارد، محدود کنیم. زبان طبیعیِ بیان
واقعیت‌های حساب~$\Lang L_A$ است. $\Lang L_A$ شامل یک !!{predicate}
دوموضعیِ~$<$، یک !!{constant}~$\Obj 0$، یک !!{function} یک‌موضعیِ~$\prime$
و دو !!{function}s دوموضعیِ~$+$ و~$\times$ است.

\begin{defn}
مجموعه‌ای از !!{sentence}s چون~$\Gamma$ یک \emph{نظریه} است اگر نسبت به
نتیجهٔ معنایی بسته باشد؛ یعنی
$\Gamma=\Setabs{!A}{\Gamma\Entails !A}$.
\end{defn}

دو راه ساده برای مشخص‌کردن نظریه‌ها وجود دارد. یکی اینکه نظریه را
مجموعهٔ !!{sentence}sی صادق در !!{structure}ی بگیریم. برای نمونه، !!{structure}ی
برای~$\Lang L_A$ در نظر بگیرید که !!{domain}اش~$\Nat$ است و همهٔ
نمادهای غیرمنطقی چنان تعبیر می‌شوند که انتظار دارید.

\begin{defn}\ollabel{def:standard-model}
\emph{مدل استاندارد حساب} !!{structure}~$\Struct N$ است که چنین تعریف
می‌شود:
\begin{enumerate}
\item $\Domain N = \Nat$
\item $\Assign{\Obj 0}{N} = 0$
\item $\Assign{\Obj \prime}{N}(n) = n + 1$ برای هر $n \in \Nat$
\item $\Assign{\Obj +}{N}(n, m) = n + m$ برای هر $n, m \in \Nat$
\item $\Assign{\Obj \times}{N}(n, m) = n\cdot m$ برای هر $n, m \in \Nat$
\item $\Assign{\Obj <}{N} = \Setabs{\tuple{n, m}}{n \in \Nat, m \in
  \Nat, n < m}$
\end{enumerate}
\end{defn}

به تفاوت میان «$\times$» و «$\cdot$» توجه کنید: $\times$ نمادی در
زبان حساب است. البته آن را چنان برگزیده‌ایم که ضرب را به یادمان
آورد، اما $\times$ عمل ضرب نیست، بلکه نماد تابعی دوموضعی است (نام
رسمی‌اش~$\Obj f^2_1$ است). در مقابل، $\cdot$ \emph{همان} تابع معمول
ضرب است. وقتی چیزی مانند~$n\cdot m$ می‌بینید، منظور حاصل‌ضرب اعداد
$n$ و~$m$ است؛ وقتی چیزی مانند~$x\times y$ می‌بینید، دربارهٔ ترمی در
زبان حساب سخن می‌گوییم. در مدل استاندارد، نماد تابعی~$\times$ به تابع
$\cdot$ روی اعداد طبیعی تعبیر می‌شود. برای جمع از~$+$ هم به‌عنوان
نماد تابعیِ زبان حساب و هم به‌عنوان تابع جمع روی اعداد طبیعی استفاده
می‌کنیم. اینجا باید از بافت بفهمید کدام معنا مقصود است.

\begin{defn}
نظریهٔ \emph{حساب صادق} مجموعهٔ !!{sentence}sیی است که در مدل استاندارد
حساب ارضا می‌شوند؛ یعنی
\[
\Th{TA} = \Setabs{!A}{\Sat{N}{!A}}.
\]
\end{defn}

$\Th{TA}$ یک نظریه است، زیرا هرگاه $\Th{TA}\Entails !A$، جملهٔ~$!A$
در هر !!{structure}ی که~$\Th{TA}$ را ارضا می‌کند ارضا می‌شود. چون
$\Sat{N}{\Th{TA}}$، داریم $\Sat{N}{!A}$ و بنابراین
$!A\in\Th{TA}$.

راه دیگر مشخص‌کردن نظریهٔ~$\Gamma$ این است که آن را مجموعهٔ
!!{sentence}sی نتیجه‌شده از مجموعه‌جمله‌ای چون~$\Gamma_0$ بگیریم. در
این صورت، $\Gamma$ «بستار»~$\Gamma_0$ نسبت به نتیجهٔ معنایی است. این
روش تنها وقتی جالب است که~$\Gamma_0$ صریحاً مشخص شده باشد، مثلاً
!!{element}sیش فهرست شده باشند. دست‌کم باید $\Gamma_0$ تصمیم‌پذیر باشد؛
یعنی آزمونی محاسبه‌پذیر وجود داشته باشد که تعیین کند یک !!a{sentence} عضو
$\Gamma_0$ هست یا نه. !!{sentence}sی~$\Gamma_0$ را \emph{اصول موضوعهٔ}
$\Gamma$ می‌نامیم و می‌گوییم $\Gamma$ به‌وسیلهٔ~$\Gamma_0$
\emph{اصل‌موضوعی شده است}.

\begin{defn}
نظریهٔ~$\Gamma$ به‌وسیلهٔ~$\Gamma_0$ \emph{اصل‌موضوعی شده است} اگر و
تنها اگر
\[
\Gamma = \Setabs{!A}{\Gamma_0 \Entails !A}.
\]
\end{defn}

\begin{defn}
نظریهٔ~$\Th{Q}$ که با جمله‌های زیر اصل‌موضوعی می‌شود «$\Th{Q}$
رابینسون» نام دارد و نظریه‌ای بسیار ساده از حساب است.
\begin{align*}
& \lforall[x][\lforall[y][(\eq[x'][y'] \lif \eq[x][y])]] \tag{$!Q_1$}\\
& \lforall[x][\eq/[\Obj 0][x']] \tag{$!Q_2$}\\
& \lforall[x][(\eq[x][\Obj 0] \lor \lexists[y][\eq[x][y']])] \tag{$!Q_3$}\\
& \lforall[x][\eq[(x + \Obj 0)][x]] \tag{$!Q_4$}\\
& \lforall[x][\lforall[y][\eq[(x + y')][(x + y)']]] \tag{$!Q_5$}\\
& \lforall[x][\eq[(x \times \Obj 0)][\Obj 0]] \tag{$!Q_6$}\\
& \lforall[x][\lforall[y][\eq[(x \times y')][((x \times y) + x)]]] \tag{$!Q_7$}\\
& \lforall[x][\lforall[y][(x < y \liff \lexists[z][\eq[(z' + x)][y]])]] \tag{$!Q_8$}
\end{align*}
مجموعهٔ !!{sentence}sی~$\{!Q_1,\dots,!Q_8\}$ اصول موضوعهٔ~$\Th{Q}$ هستند؛
پس $\Th{Q}$ شامل همهٔ !!{sentence}sیی است که از آن‌ها نتیجه می‌شوند:
\[
\Th{Q} = \Setabs{!A}{\{!Q_1, \dots, !Q_8\} \Entails !A}.
\]
\end{defn}

\begin{defn}
فرض کنید $!A(x)$ !!a{formula}ی در~$\Lang L_A$ با متغیرهای آزاد~$x$ و
$y_1$،~\dots، $y_n$ باشد. آنگاه هر !!{sentence}ای از صورت
\[
\lforall[y_1][\dots\lforall[y_n][((!A(\Obj 0) \land \lforall[x][(!A(x)
\lif !A(x'))]) \lif \lforall[x][!A(x)])]]
\]
نمونه‌ای از \emph{شِمای استقرا} است.

\emph{حساب پئانو}~$\Th{PA}$ نظریه‌ای است که با اصول موضوعهٔ~$\Th{Q}$
همراه با همهٔ نمونه‌های شِمای استقرا اصل‌موضوعی می‌شود.
\end{defn}

\begin{explain}
هر نمونهٔ شِمای استقرا در~$\Struct N$ صادق است. وقتی !!{formula}~$!A$
فقط یک !!{variable} آزاد~$x$ دارد این مطلب از همه روشن‌تر است. در این
صورت $!A(x)$ زیرمجموعه‌ای چون~$X_{!A}$ از~$\Nat$ را در~$\Struct N$
تعریف می‌کند. $X_{!A}$ مجموعهٔ همهٔ~$n\in\Nat$هایی است که وقتی
$s(x)=n$، $\Sat{N}{!A(x)}[s]$. نمونهٔ متناظر شِمای استقرا عبارت است از
\[
((!A(\Obj 0) \land \lforall[x][(!A(x) \lif !A(x'))]) \lif
  \lforall[x][!A(x)]).
\]
اگر مقدم آن در~$\Struct N$ صادق باشد، $0\in X_{!A}$ و هرگاه
$n\in X_{!A}$، آنگاه $n+1\in X_{!A}$. از $0\in X_{!A}$ به
$1\in X_{!A}$ می‌رسیم؛ از آن به~$2\in X_{!A}$ و به همین ترتیب. پس
برای هر~$n\in\Nat$، $n\in X_{!A}$. اما این یعنی
$\lforall[x][!A(x)]$ در~$\Struct N$ ارضا می‌شود.
\end{explain}

هر دو نظریهٔ~$\Th{Q}$ و~$\Th{PA}$ اصل‌موضوعی‌اند. پرسش بزرگ این است
که چقدر قوی‌اند؟ برای نمونه، آیا $\Th{PA}$ می‌تواند همهٔ حقایق
دربارهٔ~$\Nat$ را که در~$\Lang L_A$ بیان‌شدنی‌اند ثابت کند؟
مشخص‌تر، آیا اصول موضوعهٔ~$\Th{PA}$ همهٔ پرسش‌های صورت‌بندی‌شدنی در
$\Lang L_A$ را تعیین می‌کنند؟

صورت دیگر پرسش این است: آیا $\Th{PA}=\Th{TA}$؟ $\Th{TA}$ آشکارا
همهٔ حقایق دربارهٔ~$\Nat$ را اثبات می‌کند (یعنی شامل آن‌هاست) و همهٔ
پرسش‌های صورت‌بندی‌شدنی در~$\Lang L_A$ را تعیین می‌کند؛ زیرا اگر
$!A$ !!a{sentence}ای در~$\Lang L_A$ باشد، یا $\Sat{N}{!A}$ یا
$\Sat{N}{\lnot !A}$، و بنابراین یا $\Th{TA}\Entails !A$ یا
$\Th{TA}\Entails\lnot !A$. چنین نظریه‌ای را \emph{!!{complete}} بنامید.

\begin{defn}
نظریهٔ~$\Gamma$ \emph{!!{complete}} است اگر و تنها اگر به ازای هر
!!{sentence}ٔ~$!A$ در زبانش، یا $\Gamma\Entails !A$ یا
$\Gamma\Entails\lnot !A$.
\end{defn}

\begin{explain}
بنا بر قضیهٔ تمامیت، $\Gamma\Entails !A$ اگر و تنها اگر
$\Gamma\Proves !A$؛ پس $\Gamma$ کامل است اگر و تنها اگر به ازای هر
!!{sentence}ٔ~$!A$ در زبانش، یا $\Gamma\Proves !A$ یا
$\Gamma\Proves\lnot !A$.
\end{explain}

پرسش دیگری که پیش می‌آید این است: آیا روشی محاسباتی برای آزمودن عضویت
یک !!a{sentence} در~$\Th{TA}$، در~$\Th{PA}$ یا حتی فقط در~$\Th{Q}$ وجود
دارد؟ می‌توانیم با تعریف زمان !!{sentence}sبودن یک مجموعه (مثلاً
مجموعه‌ای از !!{decidable}ها) این پرسش را دقیق‌تر کنیم.

\begin{defn}
مجموعهٔ~$X$ \emph{!!{decidable}} است اگر و تنها اگر روشی محاسباتی
وجود داشته باشد که با ورودی~$x$، اگر $x\in X$ مقدار~$1$ و در غیر این
صورت~$0$ را بازگرداند.
\end{defn}

پس پرسش ما چنین می‌شود: آیا $\Th{TA}$ (یا $\Th{PA}$، یا $\Th{Q}$)
!!{decidable} است؟

پاسخ همهٔ این پرسش‌ها منفی خواهد بود. هیچ‌یک از این نظریه‌ها
تصمیم‌پذیر نیستند. اما این پدیده ویژهٔ این نظریه‌های خاص نیست. در
واقع \emph{هر} نظریه‌ای که شرایط معینی را برآورده کند مشمول همین
نتایج می‌شود. یکی از این شرایط، که $\Th{Q}$ و~$\Th{PA}$ برآورده
می‌کنند، این است که با مجموعه‌ای !!a{decidable} از اصول موضوعه
اصل‌موضوعی شده باشند.

\begin{defn}
نظریه‌ای \emph{!!{axiomatizable}} است اگر با مجموعه‌ای
!!a{decidable} از اصول موضوعه اصل‌موضوعی شده باشد.
\end{defn}

\begin{ex}
هر نظریه‌ای که با مجموعه‌ای متناهی از !!{sentence}s اصل‌موضوعی شده باشد
!!{axiomatizable} است، زیرا هر مجموعهٔ متناهی !!{decidable} است.
پس برای نمونه~$\Th{Q}$ !!{axiomatizable} است.

نظریه‌هایی مانند~$\Th{PA}$ که به‌وسیلهٔ شِما اصل‌موضوعی شده‌اند نیز
!!{axiomatizable}ند. برای آزمودن اینکه آیا $!B$ در میان اصول
موضوعهٔ~$\Th{PA}$ است، یعنی برای محاسبهٔ تابع~$\Char{X}$ که
$\Char{X}(!B)=1$ اگر $!B$ اصلی از~$\Th{PA}$ باشد و در غیر این صورت
$=0$، می‌توانیم چنین کنیم: نخست بررسی کنید آیا $!B$ یکی از اصول
موضوعهٔ~$\Th{Q}$ است. اگر هست، پاسخ «بله» و
$\Char{X}(!B)=1$ است. اگر نیست، بیازمایید که آیا نمونه‌ای از شِمای
استقراست. این کار را می‌توان منظم انجام داد؛ در این مورد شاید
ساده‌ترین راه چنین باشد: هر نمونهٔ شِمای استقرا با تعدادی سور همگانی
آغاز می‌شود و سپس زیر!!{formula}ی شرطی می‌آید. تالیِ آن شرطی
$\lforall[x][!A(x,y_1,\dots,y_n)]$ است که در آن $x$ و~$y_1$،~\dots،
$y_n$ همهٔ متغیرهای آزاد~$!A$ هستند و سورهای آغازین~$!B$ متغیرهای
$y_1$،~\dots، $y_n$ را می‌بندند. پس از استخراج این~$!A$ و بررسی
اینکه متغیرهای آزادش با متغیرهای بسته‌شده به‌وسیلهٔ سورهای همگانیِ
ابتدا و~$\lforall[x]$ مطابقت دارند، بررسی می‌کنیم که مقدم شرطی با
\[
!A(\Obj 0, y_1, \dots, y_n) \land \lforall[x][(!A(x, y_1, \dots, y_n)
\lif !A(x', y_1, \dots, y_n))]
\]
مطابقت داشته باشد. باز هم اگر مطابقت دارد، $!B$ نمونه‌ای از شِمای
استقراست و اگر ندارد، نیست.
\end{ex}

برای پاسخ‌دادن به این پرسش، و پرسش کلی‌تر دربارهٔ اینکه کدام نظریه‌ها
کامل یا تصمیم‌پذیرند، تعریف زیر نیز سودمند خواهد بود. به یاد آورید
مجموعهٔ~$X$ !!{enumerable} است اگر و تنها اگر تهی باشد یا تابعی
!!a{surjective} چون~$f\colon\Nat\to X$ وجود داشته باشد. چنین تابعی برشمردنی
از~$X$ نام دارد.

\begin{defn}
مجموعهٔ~$X$ \emph{!!{computably enumerable}} (به اختصار !!{c.e.})
نامیده می‌شود اگر و تنها اگر تهی باشد یا برشمردنی محاسبه‌پذیر داشته
باشد.
\end{defn}

افزون بر !!{axiomatizability}، شرط دیگرِ نظریه‌های مشمول قضایای
ناتمامیت این است که به‌اندازه‌ای قوی باشند که واقعیت‌های پایه دربارهٔ
توابع محاسبه‌پذیر و روابط !!{decidable} را ثابت کنند. منظورمان از
«واقعیت‌های پایه»، !!{sentence}sیی است که مقادیر توابع محاسبه‌پذیر را
برای هر آرگومان بیان می‌کنند. و «به‌اندازهٔ کافی قوی» یعنی نظریه‌های
مورد نظر این جمله‌ها را در شمار قضایای خود داشته باشند. برای نمونه،
یک تابع محاسبه‌پذیر نمونه‌وار را در نظر بگیرید: جمع. مقدار~$+$ برای
آرگومان‌های~$2$ و~$3$ برابر~$5$ است؛ یعنی $2+3=5$. جمله‌ای در زبان
حساب که این را بیان می‌کند چنین است:
$(\num{2}+\num{3})=\num{5}$. برای نمونه، $\Th{Q}$ این جمله را ثابت
می‌کند. به‌طور کلی‌تر، می‌خواهیم به ازای هر تابع محاسبه‌پذیر
$f(x_1,x_2)$، !!a{formula}ی چون~$!A_f(x_1,x_2,y)$ در~$\Lang{L_A}$ وجود
داشته باشد به‌گونه‌ای که هرگاه $f(n_1,n_2)=m$، آنگاه
$\Th{Q}\Proves !A_f(\num{n_1},\num{n_2},\num{m})$. بدین شیوه
$\Th{Q}$ ثابت می‌کند مقدار~$f$ برای آرگومان‌های~$n_1$ و~$n_2$
برابر~$m$ است. در واقع اندکی بیشتر می‌خواهیم: ثابت کند هیچ عدد دیگری
مقدار~$f$ برای آن آرگومان‌ها نیست. دربارهٔ روابط !!{decidable} نیز
همین‌گونه است. دو تعریف بعدی این را دقیق می‌کنند.

\begin{defn}
یک !!{formula}~$!A(x_1,\dots,x_k,y)$ تابع
$f\colon\Nat^k\to\Nat$ را در~$\Gamma$ \emph{!!{represents}} اگر
و تنها اگر هرگاه $f(n_1,\dots,n_k)=m$، آنگاه
\begin{enumerate}
\item $\Gamma \Proves !A(\num{n_1}, \dots, \num{n_k}, \num{m})$، و
\item $\Gamma \Proves \lforall[y](!A(\num{n_1}, \dots, \num{n_k},
y) \lif y = \num{m})$.
\end{enumerate}
\end{defn}

\begin{defn}
یک !!{formula}~$!A(x_1,\dots,x_k)$ رابطهٔ
$R\subseteq\Nat^k$ را در~$\Gamma$ \emph{!!{represents}} اگر و
تنها اگر
\begin{enumerate}
\item هرگاه $R(n_1,\dots,n_k)$،
$\Gamma\Proves !A(\num{n_1},\dots,\num{n_k})$، و
\item هرگاه $R(n_1,\dots,n_k)$ برقرار نباشد،
$\Gamma\Proves\lnot !A(\num{n_1},\dots,\num{n_k})$.
\end{enumerate}
\end{defn}

یک نظریه برای آنکه قضایای ناتمامیت بر آن اعمال شوند «به‌اندازهٔ کافی
قوی» است اگر همهٔ توابع محاسبه‌پذیر را !!{represents} و همین کار
را برای همهٔ روابط !!{decidable} نیز انجام دهد. $\Th{Q}$ و گسترش‌هایش این شرط را برآورده می‌کنند،
اما اثباتش مدتی زمان خواهد برد. این واقعیتی غیرپیش‌پاافتاده دربارهٔ
چیزهایی است که $\Th{Q}$ می‌تواند ثابت کند، و نشان‌دادنش دشوار است
زیرا $\Th{Q}$ فقط چند اصل موضوعه دارد که باید همهٔ این واقعیت‌ها را
از آن‌ها ثابت کنیم. بااین‌حال $\Th{Q}$ نظریه‌ای بسیار ضعیف است. پس
هرچند اثبات اینکه $\Th{Q}$ همهٔ توابع محاسبه‌پذیر را بازنمایی می‌کند
دشوار است، بیشتر نظریه‌های جالب از~$\Th{Q}$ قوی‌ترند؛ یعنی بیش از آن
ثابت می‌کنند. و هرچه $\Th{Q}$ ثابت کند، هر نظریهٔ قوی‌تر نیز ثابت
می‌کند؛ چون $\Th{Q}$ همهٔ توابع محاسبه‌پذیر را بازنمایی می‌کند، هر
نظریهٔ قوی‌تر نیز چنین می‌کند. پس نظریه‌های جالب بسیاری این شرطِ
قضایای ناتمامیت را برآورده می‌کنند. بنابراین کار دشوار ما نتیجه خواهد
داد، چون نشان می‌دهد قضایای ناتمامیت بر طیف وسیعی از نظریه‌ها اعمال
می‌شوند. بی‌تردید هر نظریه‌ای که بخواهد «تمام ریاضیات» را صوری کند
باید هرچه $\Th{Q}$ ثابت می‌کند ثابت کند، زیرا دست‌کم باید بتواند
نتایج محاسبات ابتدایی را بازگیرد. پس هر نظریه‌ای که نامزد نظریه‌ای
برای «تمام ریاضیات» باشد، مشمول قضایای ناتمامیت خواهد بود.


\end{document}
